% ===== §6 The Philosophy of the Body =====
\section{The Philosophy of the Body}
\label{p23:sec:body}

\noindent
The second spine holds that perception is plastic, and that the cultivation of perception is therefore a real training rather than a figure of speech. Before that claim receives its empirical body from the neuroscience of the perceiving brain (\xref{p23:sec:plasticity}), it must be given its philosophical form, for the empirical results are only intelligible as an education of perception if perception is first understood to be the kind of thing an education could reach. This section supplies that form by way of the body. It shows, from the phenomenology of embodiment and from the enactive tradition in cognitive science, that perceiving is a bodily accomplishment, a skilled coupling of a body with its world rather than a mental registration behind the eyes; and that because it is a bodily skill, it is the kind of capacity that is acquired, refined, and reshaped through the history of a body's engagements. Plasticity, in other words, is not first a neural fact awaiting philosophical interpretation. It is entailed by what perception already is, once perception is understood as embodied. The neuroscience of the next section will then be read not as the discovery that perception can be trained but as the physiological account of a trainability the philosophy of the body already requires.

\subsection{The phenomenology of the body}
\label{p23:sub:body-phenomenology}

Phenomenology distinguished, before cognitive science reached the matter, between the body one has and the body one is. The lived body, the \emph{Leib}, is not the object-body, the \emph{K\"{o}rper}, that anatomy describes and that one might survey as one surveys any thing in the world; it is the body as it is lived from within, the seat and the medium of perceiving rather than one of perception's objects. Merleau-Ponty made this lived body the ground of all perceiving. To perceive is not to receive impressions in a mind to which the body is external hardware; it is to bring to bear a bodily hold on the world, a hold organised by the body schema that arranges the perceptual field in advance of any particular perceiving within it \parencite{merleau-ponty1962}. The body schema is not a representation the perceiver consults; it is an operative, practical familiarity, a way the body already has of being geared into its surroundings, and it is this that perception draws upon. Husserl had already located the same priority in the kinaesthetic constitution of perception: the perceived thing is given as the correlate of possible movements of the perceiver's own body, so that to see the front of a thing is already to carry an implicit bodily sense of the movements that would bring its hidden sides into view \parencite{husserl2001}. Perceiving is, at its root, a matter of what the body can do and could do, a system of bodily potentialities brought to bear on what appears. This has a consequence the paper needs. If perceiving is a bodily hold, and if the hold is constituted by an acquired practical familiarity and a repertoire of kinaesthetic potentialities, then perceiving is a skill, and skills are cultivated. The lived body is not a fixed apparatus delivering a standard perceptual world; it is a history of acquired bodily competences, and its perceptual world is the correlate of that history. The everyday is where this is most fully true, since it is in the unremarkable and repeated engagements of ordinary bodily life, the ten thousand small dealings with familiar things, that the body schema is chiefly formed and reformed. What one can perceive in an ordinary room is the sediment of how one's body has learned to inhabit ordinary rooms.

\subsection{The enactive genealogy}
\label{p23:sub:body-enactive}

Cognitive science reached the same conclusion from its own side and gave it a sharper operational edge, and its convergence with phenomenology is not accidental, since the enactive programme took the reconciliation of the two as its explicit aim. The founding claim of that programme is a direct denial of the picture on which perception is the internal representation of a pre-given world. Cognition, Varela, Thompson, and Rosch proposed, is not the representation of a ready-made world by a ready-made mind but the enactment of a world and a mind together, on the basis of the history of actions a living being performs; cognition is embodied action, and perceiving is a mode of it \parencite{varela1991}. Perception brings forth a perceived world through the sensorimotor coupling of an agent with its environment rather than mirroring an environment given in advance. Two developments of this claim supply what the paper's second spine requires. The first is Gibson's ecological account, from which the enactive tradition took its start: the primary objects of perception are affordances, the possibilities for action that a thing or a place offers to a perceiver whose bodily capacities they answer to, so that to perceive is directly to perceive what the environment affords, and the perceived world is a field of action-possibilities correlative to a body's powers \parencite{gibson1979}. What one perceives is a function of what one's body can do, and a differently capable body perceives a differently afforded world. The second is the sensorimotor account of O'Regan and No\"{e}, on which perceptual content is constituted by the perceiver's practical mastery of the regularities linking movement and sensory change: to see is to exercise an implicit knowledge of how sensory input would vary with one's own action, so that what one sees is not a static internal picture but a grasp of a structured space of possible sensorimotor engagements \parencite{oregannoe2001,noe2004}. Both developments carry the same implication for cultivation, and it is the implication the whole spine was raised to secure. If perceptual content is constituted by the mastery of sensorimotor regularities, and if the perceived world is the correlate of a body's affordance-answering capacities, then perceptual capacity is precisely a mastery, an acquired competence that can be extended and refined by the enlargement of the sensorimotor repertoire it rests on. A perception is not a fixed output of a standard organism; it is the exercise of a skill whose sophistication is a function of a history of engagement, and which further engagement can deepen. The philosophy of the body thus delivers the second spine in advance of any neuroscience: because perceiving is an embodied skill, a coupling constituted by acquired sensorimotor mastery, perception is intrinsically the kind of capacity an education could cultivate. What remains is to show, in the physiology of the perceiving brain, the mechanism by which this cultivation actually occurs, and to defend it against the charge that what such a mechanism trains is not perception but mere preference. That is the work of the section that follows.
